\section{
    Revisão Bibliográfica
} \label{sec:revisao_bibliografica}

\subsection{Fundamentos da Arquitetura Web}

A Arquitetura Orientada a Serviços (SOA) organiza sistemas distribuídos na interação entre cliente, provedor e registro de serviços, normalmente no ciclo \textit{publish-find-bind} \cite{Mello2006SegurancaServicosWeb}. Nesse modelo, SOAP, WSDL e UDDI sustentam troca de mensagens, descrição de interfaces e descoberta de serviços.

Com a evolução das aplicações web, a adoção de microsserviços ampliou a granularidade e o uso de protocolos leves, como REST \cite{Pahl2016MicroservicesSurvey}. Essa transição, detalhada na Seção~\ref{sec:seguranca_arquiteturas_modernas}, aumentou a complexidade de segurança em comunicação entre serviços e exposição de APIs.

Além de ganhos de escalabilidade e manutenção, essa mudança arquitetural deslocou parte dos riscos para interfaces distribuídas e para a governança de dependências entre serviços. Em ambientes modernos, autenticação entre componentes, rastreabilidade de chamadas e consistência de políticas de acesso tornaram-se requisitos operacionais para reduzir superfície de ataque.

\subsection{Conceitos de Segurança}

A segurança da informação em aplicações web é sustentada por cinco requisitos: confidencialidade, integridade, disponibilidade, autenticidade e não-repúdio \cite{Mello2006SegurancaServicosWeb}. Esses princípios se refletem em criptografia, autenticação, autorização e mecanismos de continuidade de serviço.

Como aplicações web processam entradas externas e dados sensíveis, elas permanecem expostas a múltiplos vetores de ataque \cite{Yadav2018WebAppSecurity}. Para padronizar análise e priorização, destacam-se taxonomias como o OWASP \textit{Top} 10 \cite{OWASP2021Top10} e classificações sistemáticas de exploração de falhas \cite{Shahriar2011SecurityVulnerabilityTaxonomy}.

Essas taxonomias orientam avaliação de risco, requisitos de segurança e priorização de correções. Assim, a revisão bibliográfica descreve vulnerabilidades recorrentes e estrutura critérios para comparar estratégias de defesa.

\subsection{Vulnerabilidades e Mecanismos de Defesa}

As vulnerabilidades em aplicações web são frequentemente agrupadas por vetor de ataque. Entre as mais recorrentes estão injeções (como SQLi), \textit{Cross-Site Scripting} (XSS), falhas de controle de acesso, \textit{Cross-Site Request Forgery} (CSRF), más configurações e problemas em componentes de terceiros \cite{Halfond2006SQLInjectionClassification,Gupta2015XSSAttacksDefense,Alwan2021OWASPTopTenSurvey}.

As defesas atuam em camadas. No nível preventivo, destacam-se validação/sanitização de entradas, parametrização de consultas e \textit{tokens} anti-CSRF \cite{Shar2016WebAppDefense}. No ciclo de desenvolvimento, SAST, DAST e IAST identificam falhas em diferentes fases. Em produção, WAF e RASP ampliam a proteção em tempo de execução. A Seção~\ref{sec:mecanismos_defesa} aprofunda essas técnicas.

Nenhuma técnica isolada cobre todo o ciclo de vida da aplicação. Medidas preventivas reduzem falhas, testes especializados antecipam detecção e controles em \textit{runtime} mitigam exploração residual. Em conjunto, essas práticas sustentam uma postura contínua de segurança.
